\input{math_commands.tex}

\usepackage[hidelinks]{hyperref}
\usepackage{url}

\usepackage{adjustbox}
\usepackage[capitalise,noabbrev]{cleveref}
\usepackage{scrextend}

\usepackage{enumitem}
\usepackage{threeparttable}
\usepackage{multirow, makecell}
\usepackage{booktabs}
\usepackage{xspace}
\usepackage{wrapfig}
\usepackage{dsfont}
\usepackage[normalem]{ulem}

\usepackage{subcaption}
\usepackage{caption}
\usepackage{dashrule}
\captionsetup[figure]{font=footnotesize}

\usepackage{amsthm}
\usepackage{thmtools}
\usepackage{thm-restate}
\usepackage{mathtools}
\usepackage[makeroom]{cancel}

\declaretheoremstyle[
  spaceabove=0.5em, %
  spacebelow=0.01em,%
  headfont=\normalfont\bfseries,
  bodyfont=\normalfont,
  postheadspace=0.5em,
]{mystyle}

\declaretheoremstyle[
  spaceabove=0.5em, %
  spacebelow=0.01em,%
  headfont=\normalfont\itshape,
  bodyfont=\normalfont,
  postheadspace=0.5em,
  qed=$\square$,
]{myproofstyle}

\newtheorem{theorem}{Theorem}
\newtheorem{lemma}{Lemma}

\declaretheorem[name=Theorem,numberwithin=section,style=mystyle]{thm}
\crefname{thm}{Theorem}{Theorems}

\declaretheorem[name=Lemma,numberlike=thm,style=mystyle]{lem}
\crefname{lem}{Lemma}{Lemmas}

\declaretheorem[name=Proposition,numberlike=thm,style=mystyle]{prop}
\crefname{prop}{Proposition}{Propositions}


\declaretheorem[name=Corollary,numberlike=thm,style=mystyle]{cor}
\crefname{cor}{Corollary}{Corollaries}

\declaretheorem[name=Definition,numberlike=thm,style=mystyle]{defi}
\crefname{defi}{Definition}{Definitions}

\declaretheorem[name={Proof}, style=myproofstyle, unnumbered]{Proof}

\usepackage{tikz}
\usetikzlibrary{calc,decorations,decorations.pathmorphing}
\usetikzlibrary{positioning,fit,arrows}
\usetikzlibrary{decorations.markings}
\usetikzlibrary{shapes,shapes.geometric}
\usetikzlibrary{shadows,patterns,snakes}
\usetikzlibrary{backgrounds,decorations.pathreplacing,automata}
\usetikzlibrary{intersections}
\usetikzlibrary{angles,quotes}
\usetikzlibrary{plotmarks}
\usetikzlibrary{patterns}
\usetikzlibrary{arrows.meta}
\usetikzlibrary{shapes.misc}
\usetikzlibrary{chains}
\usepackage{siunitx}
\newcolumntype{d}[1]{S[table-format=#1]}
\usetikzlibrary{patterns.meta}

\makeatletter
\def\extractcoord#1#2#3{
	\path let \p1=(#3) in \pgfextra{
		\pgfmathsetmacro#1{\x{1}/\pgf@xx}
		\pgfmathsetmacro#2{\y{1}/\pgf@yy}
		\xdef#1{#1} \xdef#2{#2}
	};
}
\makeatother





\newcommand{\bc}[1]{\mathcal{#1}}
\newcommand{\bs}[1]{\boldsymbol{#1}}
\DeclareMathOperator*{\relu}{ReLU}
\newcommand{\RR}{\relu}

\newcommand{\I}{\mathbb{I}\xspace}
\newcommand{\cpwl}{\operatorname{CPWL}}

\usepackage{pifont}%
\newcommand{\cmark}{\ding{51}}%
\newcommand{\xmark}{\ding{55}}%

\newcommand{\starsuper}{\textsuperscript{*}}
\newcommand{\phantomstar}{\phantom{\starsuper}}

\renewcommand{\th}{\textsuperscript{th}\xspace}
\newcommand{\st}{\textsuperscript{st}\xspace}
\newcommand{\nd}{\textsuperscript{nd}\xspace}
\renewcommand{\rd}{\textsuperscript{rd}\xspace}

\newcommand{\crown}{\textsc{CROWN}\xspace}
\newcommand{\deepz}{\textsc{DeepZ}\xspace}
\newcommand{\deeppoly}{\textsc{DeepPoly}\xspace}
\newcommand{\boxd}{\textsc{IBP}\xspace}
\newcommand{\ibp}{\textsc{IBP}\xspace}
\renewcommand{\triangle}{$\Delta$\xspace}

\newcommand{\mn}{\textsc{MN}\xspace}
\newcommand{\dpo}{\textsc{DP-1}\xspace}
\newcommand{\dpO}{\textsc{DP-0}\xspace}
\newcommand{\dz}{\textsc{DZ}\xspace}

\newcommand{\diag}{\text{diag}\xspace}


\newcommand{\acnfc}{\protecting{\tikz[]{\draw[-, style=dashed](-0.1, -0.01) -- (0.22, -0.01);\node[] at (0,0) {};}\xspace}}
\newcommand{\acfc}{\protecting{\tikz[]{\draw[black,thick](-0.1, -0.01) -- (0.22, -0.01);\node[] at (0,0) {};}\xspace}}
\newcommand{\neigh}{\protecting{\tikz[]{\node[circle, fill=blue!80, inner sep=0pt, outer sep=0pt, opacity=0.5, minimum size=2.3mm]{};}\xspace}}







\DeclareMathOperator*{\din}{{\mathit{d}_{\text{in}}}}
\DeclareMathOperator*{\dout}{{\mathit{d}_{\text{out}}}}

\newcommand{\cifar}{CIFAR-10\xspace}
\newcommand{\IN}{\textsc{ImageNet}\xspace}
\newcommand{\TIN}{\textsc{TinyImageNet}\xspace}
\newcommand{\TINS}{\textsc{TinyIN}\xspace}
\newcommand{\mnist}{\textsc{MNIST}\xspace}
\newcommand{\fmnist}{\textsc{FMNIST}\xspace}

\usepackage{comment}

